\section{n4. Exact irreducible surface gluing criterion; pinch same-centre fibre and uniform coercivity barrier.}
\textit{Supports: G331, n1, n2, n3.}

Exact irreducible surface criterion and a same-centre boundary/barrier theorem. Use ENTIRE qualified G331 and n1--n3. These deductions do not assert that every nonsingleton fibre obstructs envelopes.

First, for the self-crossing SURFACE in n3 over any field of characteristic different from2, or its excellent DVR version with2 invertible, EVERY finite continuous signed obstacle f has the following exact rule. The normalization S is a semilocal DOMAIN with two regular normal SURFACE centres. Its boundary is the DISJOINT union of the two local-centre boundaries, and its C/D classes are their products by G331; disconnected boundary pieces do not require algebraic component idempotents in S. Let f\_+,f\_- be the source restrictions, H\_+,H\_- their continuous geometric envelopes, existing by the normal-surface proof in n1. The conductor quotient is the marked DVR k[z]\_(z), or V, so the only nontrivial boundary fibre is o\_+,o\_- over o. Set a=min(H\_+(o\_+),H\_-(o\_-)). Then
p(rho w)=H\_i(w) when w(u)\textless{}infinity on centre i, and p(o)=a.
A greatest geometric C or D minorant exists iff H\_+(o\_+)=H\_-(o\_-). When it exists it is continuous and belongs to C. This is an obstacle-by-obstacle NECESSARY AND SUFFICIENT condition on an irreducible nonnormal germ, not solely the explicit counterexample.

Here is the positive lower bound without pretending a fibrewise minimum is psh. If a=0, G331's normal zero dichotomy gives both H\_i=0, so zero descends and is the envelope. If a\textless{}0 and H\_i(o\_i)\textless{}0, set B\_i=(a/H\_i(o\_i))H\_i; the ratio is \textgreater{}=1 and B\_i\textless{}=H\_i. If H\_i(o\_i)=0, normal zero dichotomy gives H\_i=0 and use constant B\_i=a. These are C functions with B\_i(o\_i)=a. For n\textgreater{}=1 and an integer N\textgreater{}=max(1,n||H\_i-B\_i||), define
G\_(i,n)=max(B\_i,H\_i-w(u)/n), interpreting the second term as bottom where u is killed.
Exactly G\_(i,n)=max(B\_i,H\_i-w((u)+a\_i\textasciicircum{}N)/n). The ideal term is ACTUAL primary, with a negative effective relatively ample divisor positive on EVERY contracted integral curve; G331 addition, scaling and maxima prove membership in C, not merely pointwise nonpositivity. Equality follows since H\_i-N/n\textless{}=B\_i. The tuple has common conductor value a and is fibre constant everywhere, since all other fibres are singletons. G331 full-fibre POSITIVE descent gives a target C minorant of f. Its increasing pointwise limit equals H\_i at finite-u points and retains a at o. Increasing-limit closure is not invoked: these are lower bounds for the raw supremum. Every target D competitor pulls back below H\_i and has a common value\textless{}=both H\_i(o\_i), proving the upper bound.

If the values agree, the tuple H\_i itself descends positively to the continuous greatest envelope. If H\_i(o\_i)\textgreater{}a for some i, the finite-u locus is dense by divisorial density (u is nonzero), and continuity of H\_i makes its values along a net to o\_i converge to the larger value. The computed p is not USC. A greatest C/D minorant would equal p and hence cannot exist. This proof also applies whenever an excellent reduced local surface has finite normalization with regular SURFACE local centres, FULL infinite-conductor boundary consisting of one lift per centre over a single target point, and a nonzero conductor element u whose killed locus on each centre is just that point. The stated full-fibre and local cone hypotheses are necessary to this application; there is no arbitrary conductor minimum theorem.

Next consider the pinch family, independently computing its entire relevant normalization data. For ANY field k put
R=k[t\textasciicircum{}2,ty,y,s]\_(t\textasciicircum{}2,ty,y,s), S=k[t,y,s]\_(t,y,s).
This is a DOMAIN with presentation k[u,x,y,s]\_(u,x,y,s)/(x\textasciicircum{}2-u y\textasciicircum{}2), u=t\textasciicircum{}2,x=ty. The quadratic is irreducible since u is not a square in k(u,y,s), including in characteristic2; alternatively the displayed subring embedding and rank-two normal form prove the presentation. The integral element t satisfies t\textasciicircum{}2=u and lies in the fraction field as x/y. Thus S is finite, normal and birational, hence the normalization. Its UNIQUE marked centre is the origin: t\textasciicircum{}2,y,s vanish only there above the target maximal ideal. With A=k[u,y,s], S\_0=A direct-sum tA and R\_0=A direct-sum tyA. Membership of a+tb in R\_0 requires b divisible by y, and multiplying by t requires a divisible by y. Hence the conductor is exactly yS, which as an R-ideal is (y,x). The annihilator calculation localizes exactly. Its quotient map is
k[u,s]\_(u,s) -\textgreater{} k[t,s]\_(t,s), u -\textgreater{} t\textasciicircum{}2.
The target maximal primary gauge a=(u,x,y,s) extends to (t\textasciicircum{}2,y,s). These computations retain all source and quotient supports.

When char(k)!=2 there is an EXACT two-point FULL boundary fibre inside this ONE marked normal centre. Let alpha have support (y,t-s), beta have support (y,t+s), and give s value1 on their k[s] quotients. Both have w(aS)=min(2, infinity,1)=1, both kill the conductor, and their target restriction has support (x,y,u-s\textasciicircum{}2), with u-value2 and s-value1. The restrictions agree on EVERY element of R, since x,y are killed and t\textasciicircum{}2=s\textasciicircum{}2 in either quotient; every local unit has value0. They are distinct, since alpha(t-s)=infinity while beta(t-s)=1. The fibre has exactly these two points: modulo that target support the integral equation becomes t\textasciicircum{}2=s\textasciicircum{}2, with exactly the two primes t=±s; their quotient fraction fields are k(s), on which the given base valuation has its unique identical extension. Both centres remain the same maximal ideal (t,y,s). Thus this example is a concrete remaining same-centre conductor fibre, not an instance of n2.

Without the extra parameter s the pinch SURFACE R=k[t\textasciicircum{}2,ty,y]\_(t\textasciicircum{}2,ty,y) has conductor yS and quotient map k[u]\_(u)-\textgreater{}k[t]\_(t). Each quotient centred boundary is a singleton. Normalization restriction is therefore bijective on the FULL boundary, using G331 finite-conductor uniqueness elsewhere. At the conductor source point w(t)=1/2,w(y)=infinity for the extended gauge (t\textasciicircum{}2,y), giving target u-value1. Positive order/isometric transport and n1's arbitrary-field normal-SURFACE CE prove continuous geometric envelopes for ALL continuous signed obstacles on this irreducible nonnormal pinch surface, over EVERY field. This is a positive example with a nonprimary conductor and a separable degree-two conductor quotient when char(k)!=2. There is no assertion that its algebraic quotient map is radicial in that characteristic.

In characteristic2 the threefold quotient k[u,s]-\textgreater{}k[t,s] is finite radicial. Directly, every polynomial in t,s has square in k[u,s] without assuming k perfect, so its valuation is forced to be half the value of that square; the same holds for quotient fields and every support. The full normalization boundary map is bijective. Thus CE is equivalent to the normal-source CE, with no new arbitrary-characteristic threefold CE conclusion. The char(k)!=2 two-point fibre is not asserted in characteristic2.

There is a quantitative same-centre obstruction to the proposed METHOD, distinct from an envelope counterexample. On ANY excellent NORMAL local centre A with primary gauge a, choose a normalized centred divisorial nu. G331's normal bound supplies K\_nu\textgreater{}=1, independent of w, with nu(I)\textless{}=K\_nu w(I) for every unit-or-primary ideal and EVERY normalized w on that SAME centre. Indeed nu(f)\textless{}=C ord\_m(f), m\textasciicircum{}r contained in a, and w(m)\textgreater{}=1/r give K\_nu=C r after adjusting normalization; units are zero. Thus every geometric C/D function obeys g(w)\textless{}=g(nu)/K\_nu throughout the centre. If an obstacle F satisfies F(nu)=-D with D\textgreater{}K\_nu, EVERY minorant, and hence its raw supremum, is \textless{}=-D/K\_nu\textless{}-1 EVERYWHERE on that centre. Competitors cannot approach -1 on even a dense subset there. More specifically, if a proposed j\_n has nu(j\_n)\textgreater{}=c n, then w(j\_n)/n\textgreater{}=c/K\_nu for EVERY w, so max(-b,-1-w(j\_n)/n)\textless{}=max(-b,-1-c/K\_nu)\textless{}-1 whenever b\textgreater{}1. Powers of a nonunit cannot create a downstairs separator shallow on one region and uniformly deep at a fixed divisorial point in the same marked centre. A local element is either a unit at all centred valuations or a nonunit at all; there are no CRT units separating these lifts. This proves why n2's separate-centre construction cannot simply be reused at alpha,beta in the pinch threefold.

A useful general NORMAL/higher-dimensional deduction follows without CE. If a finite continuous signed obstacle f on an excellent normal local domain is negative anywhere, continuity and divisorial density give a centred divisorial nu with f(nu)\textless{}0. The preceding uniform K\_nu bound gives
-max(0,-min f) \textless{}= p(f)(w) \textless{}= f(nu)/K\_nu \textless{}0 for EVERY support w.
If f\textgreater{}=0, p(f)=0. Thus even the raw supremum has a uniform strict-negative bound in the nonzero case. The same conclusion holds for an excellent nonnormal germ whose finite normalization has a SINGLE marked normal local centre: pull f and every competitor to that centre. This includes the same-centre pinch test. It proves boundedness and excludes mixed-zero raw patterns; it does not prove USC, continuity, actual D attainment or greatest-minorant existence.

Finally, if source continuous envelopes are known and the source envelope of a particular pulled-back obstacle is constant on EVERY normalization fibre, G331 positive descent proves a continuous target envelope. Every target competitor pulls back below it, and its compatible descent is itself a competitor. This is a sufficient criterion in the pinch threefold with smooth algebraically closed characteristic-zero source, using n1's exact BFJ CE; necessity is not asserted there. Same-centre nonsingleton-fibre CE/appropriate envelope classification and general excellent NORMAL dimension\textgreater{}=3 CE or qualifying failure remain OPEN. Products, energy, full finite operators and literal radial identification remain separate. Artifacts n4.md and n4.tex; rerun python assemble.py; python verify\_delivery.py. The structural and coercivity claims are proofs relative to G331; finite arithmetic does not settle the residual envelope problem.

\[
 \text{greatest C/D minorant on the irreducible surface exists}
 \Longleftrightarrow H_+(o_+)=H_-(o_-).
\]
\[
 k[t^2,ty,y,s]_0\subset k[t,y,s]_0,\qquad
 \mathfrak c=yS,\quad \alpha:(y,t-s),\quad\beta:(y,t+s).
\]
\[
 g(w)\le g(\nu)/K_\nu\quad(g\in C\cup D),\qquad
 f\not\ge0\Longrightarrow p(f)(w)\le f(\nu)/K_\nu<0.
\]

